\documentclass{article}
\usepackage[margin=1in]{geometry}
\usepackage{graphicx}
\usepackage{booktabs}
\usepackage{amsmath}
\usepackage{hyperref}
\usepackage{caption}
\usepackage{float}
\usepackage{enumitem}

\title{Analysis of Popularity Bias in RAG Retrieval Systems}
\author{Amon Sisowath}
\date{February 2026}

\begin{document}

\maketitle

\begin{abstract}
As Retrieval-Augmented Generation (RAG) becomes a standard architecture for knowledge-intensive NLP tasks, understanding the biases inherent in its retrieval components is critical. This study investigates \textit{popularity bias}—the tendency of retrieval models to favor highly accessed documents over less popular ones. We isolate \textbf{intrinsic model bias} from data availability by evaluating sparse (BM25) and dense (neural embedding) retrievers on a popularity-stratified corpus of 5.9 million Wikipedia documents. Our findings reveal distinct performance characteristics across popularity deciles: while dense retrievers exhibit a "rich-get-richer" effect, favoring head entities, sparse methods offer more robust performance for long-tail queries. We provide a detailed technical analysis of the underlying mechanisms, including semantic crowding and vocabulary distribution, and discuss the implications for building fair and robust RAG systems.
\end{abstract}

\section{Introduction}

The rapid adoption of Large Language Models (LLMs) has led to the proliferation of Retrieval-Augmented Generation (RAG) systems across domains ranging from legal discovery to customer support. By grounding generation in retrieved evidence, RAG mitigates hallucinations and provides access to up-to-date knowledge. However, the reliability of a RAG system is bounded by the quality of its retriever. If the retrieval stage is biased, the downstream generation will inevitably reflect those biases.

A critical yet under-studied phenomenon in this context is \textbf{popularity bias}. In information retrieval, this bias manifests through two distinct mechanisms:

\begin{enumerate}
    \item \textbf{Coverage Bias (Extrinsic):} Arising from the natural distribution of world knowledge. Popular topics (e.g., "The Beatles") typically have orders of magnitude more written documentation than niche topics. A retriever might select popular content simply because there is statistically more of it to retrieve.
    \item \textbf{Intrinsic Model Bias:} Arising from the training process of the retrieval model itself. Even when information availability is perfectly balanced, neural models—having been pretrained on internet-scale data—may have "memorized" popular entities more effectively. This leads to higher similarity scores for popular concepts simply because they appeared more frequently during pretraining weights updates.
\end{enumerate}

\textbf{Focus of this Study:} This research specifically investigates the second mechanism: \textbf{Intrinsic Model Bias}. We aim to isolate the model's behavior from the natural distribution of data. By strictly controlling the experimental setup to test retrieval across balanced popularity deciles, we ask: \textit{In a fair setting where a model can choose any document, does it inherently prefer those that are popular?}

\section{Background}

While popularity bias has been extensively studied in Recommender Systems—where algorithms frequently suggest "blockbuster" items at the expense of niche content—there is a scarcity of research addressing this phenomenon in the specific context of RAG. As these systems move from experimental prototypes to critical production infrastructure, understanding whether they systematically under-serve long-tail queries (which often represent the most high-value, specific information needs) is essential for ensuring system robustness and fairness.

\section{Experimental Methodology}

To rigorously test for intrinsic bias, we devised an experimental design that neutralizes Coverage Bias.

\subsection{Experimental Design: Isolating Intrinsic Bias}
The core challenge in measuring model bias is separating it from data availability. If a corpus contains 1,000 documents about Entity A and only 1 about Entity B, a preference for Entity A is expected behavior (Coverage Bias).

To control for this, we stratify our corpus into \textbf{deciles of equal size} based on document popularity. In this setup, the "Head" (popular) and "Long Tail" (niche) partitions contain the exact same number of retrievable chunks. Any statistically significant performance difference between retrieving a document from Decile 10 versus Decile 1 can therefore be attributed to the \textit{intrinsic} properties of the model and its representations, rather than the volume of available data.

\subsection{Corpus Construction}
We utilized the \textbf{2019 KILT (Knowledge Intensive Language Tasks) Wikipedia snapshot}, comprising approximately 5.9 million documents.
\begin{itemize}
    \item \textbf{Popularity Annotation:} We annotated each document with page view data. To filter out short-term noise (e.g., viral news events) and capture semantic popularity, we calculated the monthly average page views over a four-year period subsequent to the snapshot.
    \item \textbf{Decile Stratification:} Documents were segmented into chunks (1000 characters, 100 overlap). These chunks were sorted by their parent page's view count and divided into ten equal-sized bins.
\end{itemize}

\subsection{Query Selection}
We aggregated queries from six diverse datasets: Natural Questions (NQ), Fever, TriviaQA, PopQA, TREX, and HotpotQA. These datasets provide \textbf{gold standard mappings}, meaning every query is paired with a specific Wikipedia page verified to contain the answer. This allows us to measure correctness (finding the specific target document) rather than just broad topical relevance.

\section{Technical Implementation}

To conduct this analysis at scale (indexing ~6M documents), we engineered a high-throughput pipeline.

\subsection{Vector Store and Search}
We utilized \textbf{Elasticsearch} as our primary vector store, leveraging its implementation of \textbf{HNSW (Hierarchical Navigable Small World)} graphs. We configured the index with an oversampling factor of 100x during Approximate Nearest Neighbor (ANN) search. This rigorous setting minimizes approximation errors, ensuring that retrieval failures are due to the vector representations themselves, not the search algorithm.

\subsection{Indexing Pipeline}
We implemented a serverless indexing architecture using **Modal**. This pipeline generated embeddings on GPUs and pushed them directly to Elasticsearch via a custom `embed-and-send` routine, bypassing local network bottlenecks. This allowed for rapid iteration when testing different embedding models.

\section{Retrieval Methods}

We evaluated two distinct paradigms to contrast their behavior:

\begin{enumerate}[label=(\alph*)]
    \item \textbf{Sparse Lexical (BM25):} The industry standard for keyword search. BM25 scores documents based on exact term matching (TF-IDF). It serves as a baseline that is theoretically "blind" to semantic popularity, only sensitive to term frequency.
    
    \item \textbf{Dense Semantic (Neural Embeddings):} We evaluated three models from the MTEB leaderboard to test if model size correlates with bias:
    \begin{itemize}
        \item \texttt{Lajavaness/bilingual-embedding-small} (384d)
        \item \texttt{intfloat/multilingual-e5-small} (384d)
        \item \texttt{intfloat/multilingual-e5-large} (1024d)
    \end{itemize}
\end{enumerate}

\section{Evaluation Metrics}

Our analysis employs a suite of metrics designed to capture different facets of retrieval performance:

\begin{itemize}
    \item \textbf{Recall@10:} The primary accuracy metric. We track this across deciles to visualize the "popularity curve" of each model.
    \item \textbf{MMR@10 (Mean Reciprocal Rank):} Assesses ranking quality. A high recall but low MRR for long-tail items would indicate that while the model \textit{finds} the correct document, it lacks the confidence to rank it highly.
    \item \textbf{Dense-BM25 Gap ($\Delta_{dense-bm25}$):} By subtracting the BM25 score from the Dense score per decile, we measure the \textit{relative advantage} of semantic retrieval. A positive slope would confirm that dense models disproportionately benefit popular entities.
    \item \textbf{Wrong Document Probability:} When the retriever fails, does it hallucinate a popular document? We measure the probability that an incorrect result comes from a higher popularity decile than the target.
    \item \textbf{Vector Space Topology ($\Delta_{sim}$):} We calculate the similarity gap between the top-1 and top-2 results. A small gap implies "semantic crowding"—many documents clustered tightly together. We investigate if popular topics suffer from this crowding more than niche ones.
\end{itemize}

\end{document}
